\documentclass[11pt,a4paper]{article}

% =============================================================================
% Provisional companion: Identification and regression results
% =============================================================================
% Standalone document collecting the per-(reform, outcome) identification
% strategies and headline regression results for the MTU15 thesis.
%
% Companion documents:
%   thesis/provisional/descriptive_facts.tex
%       Empirical foundation: who price-sets, hour-class taxonomy,
%       bid-shape distributions per technology, per-firm portfolio
%       responses, reforzada / Fase~I confounder characterisation.
%   thesis/paper/paper.tex
%       Main paper draft with the theoretical model (sec.~2) that
%       generates the per-outcome predictions tested here.
%   thesis/provisional/regression_results_legacy.tex
%       Earlier draft of the same content (DDD spec, DR-DiD, BA-vs-DDD,
%       overlap diagnostics). Tables referenced in this document are
%       still produced by the same pipeline; the legacy file is the
%       commit-history record of how the framing evolved.
%
% Status (2026-05-20): structural rewrite to a per-(reform, outcome)
% identification narrative. The identification design is iterating in
% the open --- treat the numbers here as the current state of play, not
% as final results.

\PassOptionsToPackage{table}{xcolor}
\usepackage{amsmath, amssymb, amsthm, mathtools}
\usepackage{dsfont}
\usepackage{graphicx, booktabs, subcaption, tikz, float}
\graphicspath{{../../figures/thesis/}{../../figures/working/}}
\usepackage{array, makecell, multirow}
\usepackage{xcolor}
\usepackage{longtable}
\usepackage[hidelinks]{hyperref}
\usepackage{xr-hyper}
\externaldocument{descriptive_facts}
\externaldocument{../paper/paper}
\usepackage{geometry}
\geometry{a4paper, margin=2.5cm}
\usepackage{pdflscape}
\usepackage{caption}
\captionsetup{font=small, labelfont=bf}

\usepackage{url}
\urlstyle{tt}
\Urlmuskip=0mu plus 1mu\relax
\def\UrlBreaks{\do\/\do\_\do\-\do\.\do\:\do\,}
\newcommand{\code}[1]{\nolinkurl{#1}}

\title{Provisional: identification strategies and regression results\\[3pt]\large MTU15 thesis companion (sketch of the strongest narrative)}
\author{Pablo Paramio}
\date{\today}

\begin{document}
\maketitle

\noindent\emph{\textbf{Status (2026-05-20).} Restructured from a regression-output companion into a per-(reform, outcome) identification narrative. Each block links bidding-behaviour evidence from the descriptive companion to the corresponding outcome regression, names the identification strategy, lists the threats it kills and those it leaves on the table, reports the headline coefficient, and reads the result against the model prediction. Numbers come from the existing pipeline (Stata reform-window DDD, R fPCA score regressions, Python panels); a few specifications are flagged as ``next-pass'' where the new descriptive findings demand a rerun. The design is iterating in the open and the file is meant to be read alongside the descriptive facts companion and the model in paper.tex~\S2.}

\tableofcontents

\bigskip

% =====================================================================
\section{Narrative arc and how to read this document}\label{sec:ident:arc}
% =====================================================================

\paragraph{Thesis question.} How do the MTU15 granularity reforms (MTU15-IDA, March~2025; MTU15-DA, October~2025), and their settlement companion ISP15 (December~2024), change strategic bidding behaviour and market outcomes in the Spanish electricity wholesale markets?

\paragraph{Narrative arc.} Each reform shifts the within-hour and across-stage strategy set available to firms. The descriptive companion (\code{descriptive_facts.tex}) characterises the bidding response one technology, one firm, one hour-class at a time. This document closes the loop on the outcome side: \emph{for each reform of interest, for each outcome of interest, which identification strategy isolates the reform-driven change, and is the resulting coefficient consistent with the predictions the theoretical model (paper~\S2) makes about how firms with a strategic margin in critical hours should respond?}

\paragraph{Reading by row.} The document is organised as a matrix of $(\text{reform},\,\text{outcome})$ blocks. Each block follows the same template:

\begin{enumerate}
\item \emph{What the descriptive evidence shows.} The bidding-behaviour fact that motivates studying this outcome under this reform.
\item \emph{Model prediction.} The directional or sign prediction the granularity model (paper~\S2) makes for this $(\text{reform},\,\text{outcome})$ cell.
\item \emph{Identification strategy.} The estimator (DDD, BA-with-controls, DR-DiD), the design window, the sample, and the parallel-trends or conditional-PT assumption it relies on.
\item \emph{Threats killed by the design, threats remaining.}
\item \emph{Result.} Headline coefficient with SE, supporting event-study figure, robustness columns.
\item \emph{Reading against the model.} Whether the sign and magnitude are consistent with the prediction; what the alternative readings would imply.
\end{enumerate}

\paragraph{Companion documents.} The bidding-behaviour evidence that motivates each block lives in the descriptive companion; section anchors below cite it explicitly. The model whose predictions we read against lives in \code{paper.tex}~\S2; the proof of each prediction lives there too.

\paragraph{Scope and exclusions.} This file does not re-derive descriptive evidence. It does not re-state the model. It does not chase every plausible regression --- for each outcome, exactly one identification strategy is named as the workhorse and at most two robustnesses are reported. Where the workhorse-versus-robustness divergence is informative (e.g., DR-DiD vs OLS-FE on prices, BA-with-controls vs DDD on the critical-flat spread), the divergence is noted in the reading.

% =====================================================================
\section{Model recap: what each reform should do, by outcome}\label{sec:ident:model_recap}
% =====================================================================

\paragraph{Setup (paper~\S2.1).} A finite set of generators bid into a sequential auction (day-ahead $\to$ intraday auctions $\to$ continuous intraday $\to$ balancing). Within each clock-hour, residual demand fluctuates by a tech-mix-specific amount; the cross-sectional dispersion of within-hour residual demand defines the hour-class taxonomy (descriptive companion, \S\ref{sec:prov:hour_class} therein).

\paragraph{Granularity, low and high (paper~\S\S2.2--2.3).} Under low granularity (pre-MTU15), each generator submits one price-quantity schedule per clock-hour. Under high granularity (post-MTU15), each generator submits four within-hour 15-min schedules. Firms with a price-setting margin in critical hours can in principle differentiate their within-hour bid curve to extract more rent from quarters with higher residual demand; firms that lack such a margin (flat-hour bidders, fringe price-takers) cannot.

\paragraph{Per-outcome predictions for the price-setting margin (paper~\S2.4).} Letting $\sigma^2_{\text{within}}(h)$ denote the within-hour residual-demand variance at clock-hour $h$, the model predicts:
\begin{itemize}
\item \emph{Within-hour bid-curve dissimilarity} $D_w$: rises post-MTU15-IDA for price-setting techs in critical hours; remains close to zero in flat hours and across all hours for non-price-setting techs.
\item \emph{Within-hour clearing-price standard deviation}: rises post-MTU15-DA in critical hours; remains low in flat hours.
\item \emph{Day-ahead cleared quantity} $q_1$ for the price-setting technology: ambiguous in level, but the cross-hour-class differential $q_1^{\text{crit}} - q_1^{\text{flat}}$ depends on the outside option (intraday repositioning vs system-operator post-clearing redispatch).
\item \emph{Intraday repositioning} $q_2$: positive (more intraday selling) under classical Ito--Reguant withholding; non-positive under broader disengagement where the outside option is the post-clearing channel rather than the intraday auction.
\item \emph{Critical-flat clearing-price differential}: rises (price-setters extract more rent in critical hours) or falls (fringe technologies fill the new quarter slots), depending on the cross-tech ramp constraints.
\item \emph{Imbalance volume per ISP cell}: falls post-ISP15 from the variance-reduction mechanical channel plus a behavioural channel through improved BRP forecasting.
\end{itemize}

The corollary on ramp-constrained hours (paper~\S2.4, end) amplifies the price-margin response when physical adjustment across consecutive quarters is binding: under reinforced operation (\emph{reforzada}, post-2025-04-28), the system operator imposes additional ramp restrictions, so the price-margin response should be larger post-MTU15-DA than the bare granularity-reform model would predict, with the additional component attributable to the reforzada $\times$ granularity interaction rather than to MTU15-DA alone.

% =====================================================================
\section{Reforms, outcomes, and the identification matrix}\label{sec:ident:matrix}
% =====================================================================

\paragraph{The three reforms.}
\begin{description}
\item[ISP15, 2024-12-09.] Settlement granularity moves to 15-min imbalance pricing while DA and IDA periods remain MTU60. \emph{Outcomes affected:} imbalance volumes, dual-price wedge, continuous-intraday volume.
\item[MTU15-IDA, 2025-03-19.] Intraday auctions move to MTU15 (in addition to the June-2024 6$\to$3 SIDC sessions reform). First regime where within-hour quarter bid curves exist for the IDA, hence the first regime where $D_w$ is defined. \emph{Outcomes affected:} IDA bid-shape, IDA clearing prices, intraday quarter rebalancing, continuous-intraday volume.
\item[MTU15-DA, 2025-10-01.] Day-ahead moves to MTU15. First regime where within-hour quarter bid curves exist for the DA. \emph{Outcomes affected:} DA bid-shape, DA clearing prices, $q_1$/$q_2$ allocation, post-DA gap (with reforzada confound).
\end{description}

\paragraph{Reforzada as the dominant confounder.} Post-2025-04-28 (Iberian blackout), REE has maintained an extended \emph{operaci\'on reforzada} stance: real-time PO-3.x technical-restrictions cascade plus PO-7.4 zonal reactive market (active June~2025 onward). This adds pay-as-bid redispatch volume that goes disproportionately to zonally-dominant CCGT units (descriptive companion, Part~B). The MTU15-IDA post-blackout window and the entirety of the MTU15-DA window are reforzada-on. Only the 40-day MTU15-IDA pre-blk window (2025-03-19 to 2025-04-27) is post-MTU15-IDA \emph{and} pre-reforzada, which is what makes that window the clean half of the design.

\paragraph{Identification matrix, by reform and outcome.} Table~\ref{tab:ident_matrix} lays out the per-cell strategy. The matrix tells you which cells admit a clean reform-vs-reforzada decomposition (the IDA15 design) and which carry a residual confound (the DA15 design). Within-cell details follow in \S\ref{sec:ident:isp15}--\S\ref{sec:ident:mtu15da}.

\begin{table}[H]
\centering\scriptsize
\setlength{\tabcolsep}{4pt}
\caption{Identification matrix. Each cell names the workhorse design for that $(\text{reform},\,\text{outcome})$ pair. The reforzada column says whether the reforzada confound is active in the post-window of that design; a green check on a cell means the design holds reforzada constant within the window.}
\label{tab:ident_matrix}
\begin{tabular}{@{}l l l l@{}}
\toprule
Reform & Outcome family & Identification workhorse & Reforzada in post-window? \\
\midrule
ISP15 (Dec 2024) & Imbalance volume & BA with calendar controls + ENTSO-E weather & No (pre-reforzada throughout) \\
                 & Dual-price wedge & Same                                          & No \\
                 & CI volume        & Same                                          & No \\
\midrule
MTU15-IDA (Mar 2025) & IDA price spread (crit$-$flat) & DDD on IDA-S1 + 2024 placebo, DR-DiD & No (40-day pre-blk window only) \\
                     & IDA price spread, closest session & Robustness on DR-DiD                & No \\
                     & DA price spread (spillover)       & DDD on DA price + 2024 placebo      & No \\
                     & $q_1$ (DA cleared)                & DDD on firm-(unit, date, hour)      & No \\
                     & $q_2$ (intraday repositioning)    & DDD on firm-(unit, date, hour)      & No \\
                     & $D_w$ on IDA bids                 & Pre-blk vs post-blk descriptive (not DDD; $D_w$ undefined pre-MTU15) & No \\
\midrule
MTU15-DA (Oct 2025) & DA price spread (crit$-$flat) & DDD on DA price + 2024 placebo, DR-DiD & \textbf{Yes, on/on within design} \\
                    & IDA price spread (spillover)  & Same                                   & \textbf{Yes} \\
                    & $q_1$ (DA cleared)            & DDD on firm-(unit, date, hour)         & \textbf{Yes} \\
                    & $q_2$ (intraday repositioning)& DDD on firm-(unit, date, hour)         & \textbf{Yes} \\
                    & $D_w$ on DA bids              & Pre vs post within DA15/ID15 (descriptive) & \textbf{Yes} \\
                    & Post-DA gap (Fase~I)          & Cross-regime descriptive, not DDD      & --- \\
\bottomrule
\end{tabular}
\end{table}

% =====================================================================
\section{Common identification toolkit}\label{sec:ident:toolkit}
% =====================================================================

\paragraph{Workhorse: triple-difference on the critical-flat spread with a 2024 placebo.} The two granularity reforms (MTU15-IDA, MTU15-DA) admit a within-day differencing on the (critical-hour $-$ flat-hour) spread. The 2024 placebo year on the same calendar window absorbs year-on-year seasonality in the spread itself. Let $Y_{f,u,d,h}$ denote the outcome at firm $f$, unit $u$, date $d$, clock-hour $h$, $\text{crit}_h = 1$ if $h \in \{5\text{--}9, 16\text{--}22\}$, $\text{post}_d = 1$ if $d$ is after the reform date, $\text{y25}_d = 1$ if $d$ is in the 2025 treatment window. Restricting the sample to $h \in \{\text{crit},\,\text{flat}\}$:
\begin{multline}
Y_{f,u,d,h} = \alpha + \beta_1 \text{crit}_h + \beta_2 \text{post}_d + \beta_3 \text{y25}_d
+ \beta_4 (\text{crit} \times \text{post}) + \beta_5 (\text{crit} \times \text{y25}) \\
+ \beta_6 (\text{post} \times \text{y25})
+ \boxed{\beta_7 (\text{crit} \times \text{post} \times \text{y25})}
+ \gamma_f + \mu_u + \delta_{\text{DOW}(d)} + \varepsilon_{f,u,d,h}.
\label{eq:ddd}
\end{multline}
$\beta_7$ is the headline coefficient. SE clustered at the date level. The estimand and a careful reading of what it captures (absolute critical-hour effect vs differential between two treated arms) are discussed in \S\ref{sec:ident:reading_DDD}.

\paragraph{Reform windows.} The two granularity-reform windows hold reforzada constant by construction within their pre/post sub-windows:
\begin{itemize}
\item \emph{IDA15 design}: pre = 2024-12-09 to 2025-03-18 (ISP15-win, reforzada off); post = 2025-03-19 to 2025-04-27 (40-day pre-blk, reforzada off); 2024 placebo = same calendar dates one year earlier.
\item \emph{DA15 design}: pre = 2025-04-28 to 2025-09-30 (post-blk, reforzada on); post = 2025-10-01 to 2026-02-13 (DA15/ID15, reforzada on); 2024 placebo = same calendar dates one year earlier (reforzada off in 2024).
\end{itemize}
Each design holds reforzada \emph{categorically} constant within the 2025 window, but the DA15 design admits a within-design reforzada-intensity drift (Fase~I cost rose roughly $2.8\times$ between the DA15 pre and post sub-windows; descriptive companion, \S\ref{sec:prov:full_cascade:perday}). The reforzada-intensity drift is the most consequential threat to the DA15 design and is addressed by the placebo + outcome-specific BA cross-check in \S\ref{sec:ident:mtu15da}.

\paragraph{Estimators.} For each $(\text{window},\,\text{outcome})$ cell we report two estimators of the same DDD parameter:
\begin{itemize}
\item \emph{OLS-FE} on the period-level panel: \texttt{reghdfe} with clock-hour, day-of-week, and (optionally) calendar-month or date fixed effects, SE clustered at date.
\item \emph{DR-DiD} (Sant'Anna and Zhao, 2020) on the date-level spread panel: two 2$\times$2 doubly-robust DiDs, one on each hour-class, with $X = (\text{i.dow}, \text{i.month})$ covariates; the DDD is the difference of the two ATTs with joint SE from the recentered influence functions. The implementation is \texttt{drdid, drimp}; details and the closest-to-delivery IDA robustness are in \S\ref{sec:ident:price_robustness}.
\end{itemize}
Without covariates the two estimators coincide. With $X = (\text{dow}, \text{month})$ the DR-DiD reweights for residual calendar imbalance (year-on-year drift in day-of-week alignment and DST boundaries); the discrepancy between the two point estimates is the IPW contribution.

\paragraph{Hour-class taxonomy.} Critical hours: $h \in \{5,6,7,8,9,16,17,18,19,20,21,22\}$. Flat hours: $h \in \{1,2,3\}$. Midday hours $h \in \{11,12,13,14\}$ are excluded from the control set because solar dominates the price-setting technology there, mixing the granularity and the operational/solar channels. The taxonomy is anchored to within-hour residual-demand variance from ENTSO-E A65+A75 data and is described in detail in the descriptive companion (\S\ref{sec:prov:hour_class}).

\paragraph{Event study (visual pre-trends).} For each $(\text{reform},\,\text{outcome})$ cell we plot the monthly within-day differential $\bar Y^{\text{crit}}_m - \bar Y^{\text{flat}}_m$ separately for the 2025 treatment year and the 2024 placebo year, anchored to $\tau = $ months from the reform date and normalised by subtracting the $\tau = -1$ value so both lines start at zero. The test: for $\tau < 0$ the two lines should overlap (parallel pre-trends); at $\tau \ge 0$ the persistent vertical separation between them is the treatment effect.

\paragraph{Doubly-robust outcome and propensity models.} The DR-DiD propensity model is a logit of $\text{y25}$ on $(\text{i.dow},\,\text{i.month})$; the outcome model is WLS on the same covariates. Overlap is essentially balanced under calendar-matched windows (median $\hat p(X) \approx 0.50$, no observations below $0.10$ or above $0.90$); the IPW component is not driven by tail behaviour of the propensity score. Per-cell overlap histograms are produced as working figures.

% =====================================================================
\section{Reading conventions and the absolute-vs-differential question}\label{sec:ident:reading_DDD}
% =====================================================================

\paragraph{What the DDD identifies.} The within-day differencing on (critical $-$ flat) absorbs day-level common shocks; the 2024 placebo absorbs year-on-year seasonal variation in the critical-flat spread itself. The identified parameter is the (critical $-$ flat) differential reform response, net of the calendar baseline. Whether that differential equals the absolute critical-hour reform effect depends on whether flat hours are themselves exposed to the reform.

\paragraph{New evidence (descriptive companion, $D_w$ ridges).} Across the three MTU15-IDA regimes, CCGT within-hour $D_w > 0$ rate in flat hours moves $10.4\% \to 30.2\% \to 50.6\%$, against $21.9\% \to 44.7\% \to 70.7\%$ in critical hours. \emph{Flat hours are not MTU15-untreated under DA15/ID15.} The DDD on the DA15 design therefore identifies the (critical $-$ flat) differential between \emph{two} treated arms --- a lower bound on the absolute critical-hour effect (paper~\S3.6 / OJO comment), not its approximation.

This is the key change in framing relative to the legacy version of this document. The DA15 coefficient is still a clean estimate, but its \emph{interpretation} downgrades from ``absolute critical-hour effect'' to ``critical-flat differential''. The level of the absolute effect on critical hours is bounded by combining the DDD with an absolute-effect calculation from the BA-with-controls spec (\S\ref{sec:ident:mtu15da}).

\paragraph{Pivotal-status heterogeneity, revised.} The earlier reading was that within-hour granularity exploitation was concentrated at Naturgy CCGT; the extended $D_w$ data (descriptive companion, \S\ref{sec:prov:dw_tables}) shows that under DA15/ID15 \emph{all four} pivotal CCGT firms have $D_w > 0$ in 59--79\% of critical-hour cells. The pivotal-vs-non-pivotal contrast in the $q_1$/$q_2$ DDD remains a heterogeneity decomposition --- not a fourth-difference identification --- but the granularity-exploitation channel is no longer single-firm.

\paragraph{Naming conventions.} We use ``price-setter'' rather than ``pivotal'' throughout (the EUPHEMIA partial-acceptance concept is the relevant primitive); the older ``pivotal'' label is preserved only in citations to the legacy regression output, where the variable name still appears. We do not use ``treatment'' or ``placebo'' to label \emph{firms} --- firms are subjects partitioned by economic identity (Big-4 dominant operators vs other vertically-integrated operators), not assigned to a treatment arm. ``Treatment'' refers exclusively to the time dimension (post-reform vs pre-reform) and the within-day dimension (critical hours vs flat hours).

% =====================================================================
\section{ISP15 (Dec 2024): identification by outcome}\label{sec:ident:isp15}
% =====================================================================

ISP15 changed the granularity of imbalance settlement, not of bidding. Its identifying contrast is BA-with-controls (the 2024 placebo cannot help here because the reform happened in 2024). The outcomes affected are on the balancing side; the relevant within-day differencing is across (critical, flat) hour-classes when residual-demand variance feeds imbalance volume.

\paragraph{ISP15 / imbalance volume.}
\begin{itemize}
\item \emph{Descriptive evidence.} Per-(BRP, ISP) imbalance MWh fell roughly 60\% per cell after ISP15 (\S\ref{sec:prov:full_cascade:efficiency}); within-hour variance decomposition attributes $\sim 50\%$ to the mechanical variance-reduction channel and $\sim 10$~pp to a behavioural improvement.
\item \emph{Model prediction.} Imbalance volume should fall: under MTU15 the within-hour variance of programmed-minus-actual is mechanically smaller because the four 15-min slots track residual demand more closely; the behavioural component reflects BRPs investing in finer-grained forecasting.
\item \emph{Identification.} BA-with-controls on monthly per-cell imbalance volume; calendar-month FE absorb seasonality, weather covariates (ENTSO-E A65 load, A75 VRE) absorb the supply-side variance shifters. PT here is the BA assumption that the pre-reform monthly cell volume would have continued at the pre-reform level absent the reform.
\item \emph{Result.} Headline: $-60\%$ raw, decomposed as $-50\%$ mechanical + $-10$~pp behavioural after VRE controls. The decomposition method and full table are in the descriptive companion (\S\ref{sec:prov:full_cascade:efficiency}).
\item \emph{Reading vs model.} Consistent with the model's mechanical-plus-behavioural decomposition. The behavioural residual is the BRP forecasting margin, which the model treats as exogenous and which the data confirms is a small but real component.
\end{itemize}

\paragraph{ISP15 / dual-price wedge.}
\begin{itemize}
\item \emph{Descriptive evidence.} The up/down dual-price wedge widened from $\approx 40/85$ to $\approx 70/110$ EUR/MWh; per-regime price table in \S\ref{sec:prov:full_cascade:prices}.
\item \emph{Model prediction.} The wedge widens if granular settlement increases the cost of being short relative to being long --- a function of the system-operator's pricing rule rather than of firm strategy.
\item \emph{Identification.} BA-with-controls on the monthly mean dual prices, ENTSO-E weather as covariates. Not a DDD because the reform affects all hours symmetrically.
\item \emph{Result.} The +30/+25 EUR/MWh wedge widening is robust to the inclusion of VRE controls and gas-price controls (\S\ref{sec:prov:full_cascade:prices}).
\item \emph{Reading vs model.} The wedge widening was anticipated by the design of the dual-price imbalance rule; the magnitude is informative for the welfare calculation in the cost-cascade section.
\end{itemize}

\paragraph{ISP15 / continuous-intraday volume.}
\begin{itemize}
\item \emph{Descriptive evidence.} CI volume fell $-15\%$ SA at ISP15 (\S\ref{sec:prov:ci_seasonality_adjusted}); a further fall to half its baseline followed under MTU15-IDA.
\item \emph{Model prediction.} CI provides intra-day forecast correction. Improved BRP forecasting under ISP15 mechanically reduces the demand for CI repositioning. The model does not generate a sharp directional prediction here.
\item \emph{Identification.} Seasonality-adjusted regime mean comparison from the FWL panel (descriptive companion, \S\ref{sec:bid_internal:calendar_discipline}); no DDD because the within-day contrast is not the active margin (CI volumes are aggregated to the system level for this outcome).
\item \emph{Result.} See \S\ref{sec:prov:ci_seasonality_adjusted}, Table~\ref{tab:ci_volume_seasonality_adjusted}.
\item \emph{Reading vs model.} Consistent with the BRP-forecasting interpretation of ISP15; the further fall under MTU15-IDA is consistent with the granular IDA absorbing within-hour scarcity that previously needed CI correction.
\end{itemize}

% =====================================================================
\section{MTU15-IDA (Mar 2025): identification by outcome}\label{sec:ident:mtu15ida}
% =====================================================================

MTU15-IDA is the cleaner half of the granularity-reform identification. The post-window (2025-03-19 to 2025-04-27) is pre-reforzada throughout, so the design holds reforzada at \emph{off} within both pre and post sub-windows. The 2024 placebo on the same calendar dates absorbs the YoY seasonality of the (critical $-$ flat) spread.

\paragraph{MTU15-IDA / IDA price spread (critical $-$ flat).}
\begin{itemize}
\item \emph{Descriptive evidence.} The within-hour quarter-dissimilarity $D_w$ first appears under this reform (60-min IDA periods pre-MTU15-IDA give no within-hour quarters). CCGT critical $D_w > 0$ jumps from $21.9\%$ pre-blackout to $44.7\%$ post-blackout (Fig.~\ref{fig:dw_3d_hourly_ccgt} in the descriptive companion); the price-setting margin is exercised post-MTU15-IDA. The flip in the critical-flat IDA in-band share ordering (critical below flat at 3-sess, critical above flat at DA15/ID15; SA-robust per the deleted Table~14 reading) is the cleanest cross-class signal across the reforms.
\item \emph{Model prediction.} If price-setters in critical hours use the new within-hour pricing dimension, the IDA critical-flat clearing-price spread should rise; if non-pivotal generators (renewables, fast-ramp hydro) fill the new 15-min slots competitively, the spread should fall.
\item \emph{Identification.} DDD on the IDA-S1 clearing-price spread per Eq.~\ref{eq:ddd}, 2024 placebo on the matching calendar windows. The IDA-S1 choice is robustness-checked against the closest-to-delivery session (\S\ref{sec:ident:price_robustness}) because June~2024 SIDC reform changed session 1's clearing horizon.
\item \emph{Threats killed.} Day-level cost shocks (within-day differencing), year-on-year seasonal swing of the spread (2024 placebo), reforzada (off in both pre and post), and asymmetric TTF gas confound (descriptive: CCGT price-setting share 9.1\% critical vs 9.0\% flat, no asymmetry).
\item \emph{Threats remaining.} (i) The post-window is only 40 days; the SE on the DDD coefficient reflects this. (ii) IDA-S1 horizon changed at the June~2024 SIDC reform inside the placebo year; partially absorbed by the closest-to-delivery robustness.
\item \emph{Result.} DR-DiD point estimate $+9.43$ EUR/MWh (SE $2.95$, $p < 0.01$); OLS-FE $\beta_7 = +0.50$ (SE $3.69$, n.s.); see Table~\ref{tab:reform_ddd_summary} below. The closest-to-delivery robustness shrinks the DR-DiD to $+6.26$ ($p = 0.064$): direction and order of magnitude robust, statistical significance is at the margin. We treat the headline $+9.43$ as an upper bound on the true differential.
\item \emph{Reading vs model.} The spread widens, consistent with the price-setter-exploits-granularity prediction. The asymmetric across-hour-class response (critical rises more than flat falls; descriptive companion, \S\ref{sec:bid_internal:calendar_discipline}) is also consistent with the model. The reduced-form result alone does not separate the price-setter channel from the fringe-fill channel; the per-firm $q_1$/$q_2$ regressions below pin which compositional story is consistent with the data.
\end{itemize}

\paragraph{MTU15-IDA / DA price spread (spillover).}
\begin{itemize}
\item \emph{Descriptive evidence.} See \S\ref{sec:prov:steepness} for the cross-market pattern of in-band share. DA bidding may respond to a change in the value of unloaded MW being repositioned later.
\item \emph{Model prediction.} Direct effect is zero (DA is MTU60 throughout the IDA15 window). Spillover would arise if DA bidders anticipate the new IDA repositioning value.
\item \emph{Identification.} Same DDD as the IDA price cell, applied to the DA clearing price spread.
\item \emph{Result.} DR-DiD point estimate $+6.28$ EUR/MWh (SE $3.01$, $p < 0.05$); OLS-FE $\beta_7 = -2.24$ (n.s.). The DA-side number is unaffected by the IDA-session robustness.
\item \emph{Reading vs model.} Suggestive of forward-looking DA bidding under anticipation of MTU15-IDA, but the OLS-FE/DR-DiD divergence is large enough that we read this as an open empirical question rather than a firm finding.
\end{itemize}

\paragraph{MTU15-IDA / firm-level $q_1$ (DA-cleared MWh).}
\begin{itemize}
\item \emph{Descriptive evidence.} The Big-4 CCGT $q_1$ collapse begins in earnest after the blackout, not at MTU15-IDA itself; descriptive companion, \S\ref{sec:prov:q1_drop}.
\item \emph{Model prediction.} If price-setters use the new IDA pricing dimension to withhold from DA and reposition in IDA at higher per-quarter prices, $q_1$ should fall (Ito-Reguant signature).
\item \emph{Identification.} DDD on $q_1$ at the (firm, unit, date, hour) panel with unit + clock-hour + DOW FE; firm-cluster SE.
\item \emph{Result.} $\beta_7 = -9.96^{***}$ MWh/unit-hour (SE $3.10$); pivotal-status contrast $\delta_{\text{piv}} = -4.01$ (n.s., $G \approx 9$ wild-bootstrap SE flagged as next pass). All CCGT firms move in the same direction, no market-power moderation.
\item \emph{Reading vs model.} Compatible with broader disengagement (descriptive companion, \S\ref{sec:prov:q1_drop:components}) --- DA $q_1$ falls, but the IDA $q_2$ does not rise (next paragraph); the outside option is not the IDA, it is the post-clearing redispatch channel.
\end{itemize}

\paragraph{MTU15-IDA / firm-level $q_2$ (intraday repositioning).}
\begin{itemize}
\item \emph{Descriptive evidence.} Intraday programmed MWh net of DA does not rise; the post-DA gap (PHF $-$ PDBF), which captures the Fase~I redispatch channel, does rise (descriptive companion, \S\ref{sec:prov:phf_minus_pdbf}).
\item \emph{Model prediction.} $\beta_7 > 0$ under classical Ito-Reguant withholding; $\beta_7 \le 0$ under broader disengagement.
\item \emph{Identification.} DDD on signed intraday programmed MWh at the (firm, unit, date, hour) panel.
\item \emph{Result.} $\beta_7 = -1.70^{**}$ MWh/unit-hour (SE $0.85$), \emph{negative}. The pivotal-status contrast $\delta_{\text{piv}} = +11.10$ is large but with wide SE ($10.81$), uninformative.
\item \emph{Reading vs model.} The $(q_1 < 0,\,q_2 < 0)$ quadrant is the broader-disengagement quadrant, not Ito-Reguant withholding. The outside option for CCGT is the post-clearing channel.
\end{itemize}

% =====================================================================
\section{MTU15-DA (Oct 2025): identification by outcome, with the reforzada confound}\label{sec:ident:mtu15da}
% =====================================================================

MTU15-DA is the bruised half of the granularity identification. The 2025 pre/post windows are both reforzada-on, but reforzada itself intensified within the design: Fase~I cost roughly tripled between the DA15 pre and post sub-windows. The DDD attributes that intensification to MTU15-DA unless the reforzada drift is held separately constant. The 2024 placebo does not help with this specific confound because 2024 has no reforzada at all.

\paragraph{MTU15-DA / DA price spread (critical $-$ flat).}
\begin{itemize}
\item \emph{Descriptive evidence.} CCGT critical $D_w > 0$ jumps from $44.7\%$ (post-blackout) to $70.7\%$ (DA15/ID15); flat hours also rise to $50.6\%$ (Fig.~\ref{fig:dw_3d_hourly_ccgt}). All four pivotal CCGT firms are within-hour-active under DA15/ID15.
\item \emph{Model prediction.} The critical-flat spread direction depends on (a) whether price-setters exploit the new pricing dimension to raise critical prices and (b) whether fringe technologies (renewables, fast-ramp hydro) fill the new critical-hour quarter slots and compete the spread down.
\item \emph{Identification.} DDD on the DA clearing-price spread per Eq.~\ref{eq:ddd}. DA15 window: 2025-04-28 to 2026-02-13.
\item \emph{Threats killed.} Day-level cost shocks (within-day differencing), year-on-year seasonal swing (2024 placebo). The reforzada \emph{categorical} confound is held off (on/on), but the within-design reforzada-intensity drift is not.
\item \emph{Threats remaining.} (i) Reforzada-intensity drift (the dominant confound; addressed below by a separate Fase~I cost trend regression). (ii) ``Flat hours are MTU15-untreated'' assumption is broken; the DDD coefficient now identifies the differential between two treated arms, not the absolute critical-hour effect.
\item \emph{Result.} DR-DiD $-5.38^{**}$ EUR/MWh (SE $2.36$); OLS-FE $\beta_7 = -0.29$ (n.s.). The signs are consistent in direction; the magnitude divergence is the IPW reweighting on calendar covariates.
\item \emph{Reading vs model.} Spread compresses, consistent with the fringe-fill channel. The (critical price falls more than flat price) decomposition into 2$\times$2 DiDs ($-8.34$ critical vs $-2.96$ flat under DR-DiD) reinforces the fringe-fill reading: it is the critical-hour price that falls, not the flat-hour price that rises. The interpretation as the absolute critical-hour effect is downgraded as noted in \S\ref{sec:ident:reading_DDD}; the DA15 DDD is a critical-vs-flat differential of moderate magnitude.
\item \emph{Reforzada-intensity check (next-pass).} A separate BA on Fase~I cost / TR cost between the DA15 pre and post sub-windows (currently in the cost-cascade section of the descriptive companion, \S\ref{sec:prov:full_cascade:perday}) gives a bound on the reforzada-intensity confound. The cleaner version of this check pulls reforzada-intensity into the DDD as a continuous covariate, replicating the Sant'Anna-Zhao $X$ argument for a continuous reforzada index.
\end{itemize}

\paragraph{MTU15-DA / IDA price spread (spillover).}
\begin{itemize}
\item \emph{Descriptive evidence.} Continuous-intraday volume partially recovered post-MTU15-DA; per-tech CI evolution favours RES (Wind $+36\%$, Solar PV $+108\%$ SA vs MTU15-IDA post-blk; \S\ref{sec:prov:ci_seasonality_adjusted}).
\item \emph{Identification.} Same DDD applied to IDA-S1 price.
\item \emph{Result.} DR-DiD $-5.36^{**}$ EUR/MWh (SE $2.30$); same direction and near-identical magnitude as the DA cell. Closest-to-delivery robustness drops this to $+1.21$ (n.s.); the IDA-spillover number is not robust at all.
\item \emph{Reading vs model.} The DA-IDA cross-market consistency in the DA15 design is suggestive of a real spillover but is fragile to the session-choice robustness. We read this cell as a flag, not a finding.
\end{itemize}

\paragraph{MTU15-DA / firm-level $q_1$ (DA-cleared MWh).}
\begin{itemize}
\item \emph{Descriptive evidence.} The CCGT $q_1$ collapse is concentrated at the DA15 window: $-83\%$ cal-matched (Dec 2024 vs Dec 2025; \S\ref{sec:prov:q1_drop}). The dispatch substitutes onto post-DA Fase~I redispatch (\S\ref{sec:prov:phf_minus_pdbf}).
\item \emph{Model prediction.} $\beta_7 < 0$ under any of the three CCGT-disengagement stories the model permits.
\item \emph{Identification.} DDD on $q_1$ at the (firm, unit, date, hour) panel.
\item \emph{Result.} Pooled $\beta_7 = -16.27^{***}$ MWh/unit-hour (SE $2.69$); pivotal-status contrast $\delta_{\text{piv}} = -18.65^{**}$ (SE $7.35$): the entire pooled coefficient is concentrated in the pivotal firms; non-pivotal CCGT firms in the same window show no measurable change.
\item \emph{Reading vs model.} The pivotal-only $q_1$ collapse, combined with the post-DA gap rise of similar order of magnitude (\S\ref{sec:prov:phf_minus_pdbf}, $\sim 1{,}044$ GWh/mo post-blackout), is consistent with the broader-disengagement narrative: pivotal CCGT firms reduce DA participation because the outside option (Fase~I pay-as-bid redispatch on zonally-dominant units) becomes more attractive. The mechanism's economic content is the markup over marginal cost on the Fase~I leg, documented in the descriptive companion at $\$~30$--$60\%$ over CCGT SRMC.
\end{itemize}

\paragraph{MTU15-DA / firm-level $q_2$ (intraday repositioning).}
\begin{itemize}
\item \emph{Descriptive evidence.} Intraday programmed volume net of DA does not rise; the substitution channel is post-DA, not IDA.
\item \emph{Identification.} DDD on signed intraday programmed MWh.
\item \emph{Result.} Pooled $\beta_7 = -1.27^{**}$ MWh/unit-hour; pivotal contrast $\delta_{\text{piv}} = +18.78$ with very wide SE ($15.82$) --- not identified.
\item \emph{Reading vs model.} Same broader-disengagement quadrant as MTU15-IDA. The model's outside-option-is-IDA prediction is rejected for CCGT; the data is consistent with outside-option-is-Fase~I-redispatch.
\end{itemize}

\paragraph{MTU15-DA / post-DA gap (Fase~I channel).}
\begin{itemize}
\item \emph{Descriptive evidence.} Per-(firm, tech, date, hour) post-DA gap SA panel; per-tech hour-conditional ridges in Figs.~\ref{fig:post_da_gap_diurnal}--\ref{fig:post_da_gap_3d_hourly_hp} (descriptive companion).
\item \emph{Model prediction.} The model treats reforzada as a parallel mechanism, not a granularity outcome. The post-DA gap response to MTU15-DA is therefore expected to be small in the absence of an interaction.
\item \emph{Identification.} Not a DDD --- the post-DA gap is a Fase~I outcome that responds primarily to reforzada, not to granularity. We characterise the cross-regime distribution descriptively (the hour-conditional ridges show the reform-window mass migration concentrated at the two critical hours).
\item \emph{Reading vs model.} The reforzada intensification within the DA15 window is the substantive driver. The MTU15-DA reform itself probably has a small direct effect on this outcome.
\end{itemize}

% =====================================================================
\section{Headline DDD table and event study}\label{sec:ident:headline}
% =====================================================================

\paragraph{Headline DDD on price levels, all four (window, market) cells.}

\begin{table}[H]
\centering\footnotesize
\setlength{\tabcolsep}{4pt}
\caption{Headline DDD on prices in levels (EUR/MWh), all four (reform window, market) cells. DR-DiD on the date-level spread panel; OLS-FE $\beta_7$ on the period-level panel. SE in parentheses (clustered date). Stars: $^{*}~0.10,~^{**}~0.05,~^{***}~0.01$. Channel: direct = the market hit by the reform; spillover = cross-market test.}
\label{tab:reform_ddd_summary}
\begin{tabular}{llcccc}
\toprule
& & \multicolumn{3}{c}{DR-DiD on levels} & OLS-FE \\
\cmidrule(lr){3-5}
Window & Outcome & ATT$_{\text{crit}}$ & ATT$_{\text{flat}}$ & DDD & $\beta_7$ \\
\midrule
DA15  & DA price        & $-8.34^{*}$        & $-2.96$        & $-5.38^{**}$ & $-0.29$ \\
DA15  & IDA-S1 price    & $-10.38^{**}$      & $-5.02$        & $-5.36^{**}$ & $-2.67$ \\
IDA15 & DA price        & $-25.02^{***}$     & $-31.30^{***}$ & $+6.28^{**}$ & $-2.24$ \\
IDA15 & IDA-S1 price    & $-25.40^{***}$     & $-34.82^{***}$ & $+9.43^{***}$ & $+0.50$ \\
\bottomrule
\end{tabular}
\end{table}

\paragraph{Headline DDD on firm-level $q_1$ / $q_2$.}

\begin{table}[H]
\centering\footnotesize
\setlength{\tabcolsep}{4pt}
\caption{Pooled DDD ($\beta_7$) and pivotal-status heterogeneity contrast ($\delta_{\text{piv}}$) on firm-level $q_1$ and $q_2$, MWh per unit-hour. $\delta_{\text{piv}} > 0$ means pivotal-firm response is larger in algebraic terms. SE clustered at date (pooled) and at firm (decomposed; $G \approx 9$). Stars: $^{*}~0.10,~^{**}~0.05,~^{***}~0.01$.}
\label{tab:reform_firm_summary}
\begin{tabular}{llcc|cc}
\toprule
& & \multicolumn{2}{c}{Pooled DDD ($\beta_7$)} & \multicolumn{2}{c}{$\delta_{\text{piv}}$ (heterogeneity)} \\
\cmidrule(lr){3-4}\cmidrule(lr){5-6}
Outcome & Window & $\beta_7$ & SE & $\delta_{\text{piv}}$ & SE \\
\midrule
$q_1$ & DA15  & $-16.27^{***}$ & (2.69) & $-18.65^{**}$  & (7.35) \\
$q_1$ & IDA15 & $-9.96^{***}$  & (3.10) & $-4.01$        & (6.55) \\
$q_2$ & DA15  & $-1.27^{**}$   & (0.64) & $+18.78$       & (15.82) \\
$q_2$ & IDA15 & $-1.70^{**}$   & (0.85) & $+11.10$       & (10.81) \\
\bottomrule
\end{tabular}
\end{table}

\paragraph{Event study.}
\begin{figure}[H]
\centering
\includegraphics[width=\textwidth]{fig_event_study_ddd.pdf}
\caption{Event-study, monthly (critical $-$ flat) within-day differential, normalised to $\tau = -1$. Rows = outcomes (DA price, IDA price, $q_1$, $q_2$); columns = reforms (IDA15, DA15). Red = 2025 treatment, blue = 2024 placebo. Pre/post regions lightly shaded.}
\label{fig:event_study}
\end{figure}

\textit{Pre-trend reading.} DA15: pre-period overlap is reasonable for prices; $q_1$ shows the 2025 series sitting below 2024 already before $\tau = 0$, signalling that the reforzada channel pre-bends the differential before MTU15-DA fires --- a partial PT violation for $q_1$ on this window. IDA15: pre-period overlap reasonable for prices; post-period too short to read visually. Wild-cluster bootstrap on the headline DDD is flagged as the next pass.

\paragraph{BA-with-controls vs DDD: how much of the apparent reform effect is seasonal?}
\begin{table}[H]\centering\footnotesize
\setlength{\tabcolsep}{4pt}
\caption{Naive before-after (BA) vs DDD on the within-day (crit$-$flat) spread, EUR/MWh, all four (window, market) cells. BA: within-year (post $-$ pre) on 2025 spread, OLS with \code{i.dow + i.month} calendar controls and date-clustered SE. Placebo BA: same spec on 2024 spread. DDD: 2$\times$2 DiD with both years and same calendar controls (the headline OLS-FE coefficient on $\beta_{\text{post}\times\text{y25}}$). \emph{Seasonal bias} = BA $-$ DDD = what an advisor's naive within-year BA-with-controls attributes to the reform but is really seasonality (identified by the 2024 placebo). Standard errors in parentheses. Stars: $^{*}~0.10,~^{**}~0.05,~^{***}~0.01$.}
\label{tab:ba_vs_ddd}
\begin{tabular}{llcccc}
\toprule
Window & Outcome & BA 2025 & BA 2024 (placebo) & DDD & Seasonal bias \\
       &         & (a)     & (b)               & (c) & (a) $-$ (c)   \\
\midrule
DA15 & DA price &   25.83*** \scriptsize{( 3.72)} &   26.99*** \scriptsize{( 4.34)} &   -0.17 \scriptsize{( 2.41)} &   26.00 \\
DA15 & IDA-S1 price &   25.90*** \scriptsize{( 3.79)} &   26.42*** \scriptsize{( 3.99)} &   -2.24 \scriptsize{( 2.33)} &   28.15 \\
IDA15 & DA price &    2.24 \scriptsize{( 5.30)} &   -8.49*** \scriptsize{( 3.21)} &   -1.74 \scriptsize{( 3.35)} &    3.98 \\
IDA15 & IDA-S1 price &    6.07 \scriptsize{( 4.91)} &   -7.05** \scriptsize{( 3.25)} &    0.74 \scriptsize{( 3.31)} &    5.33 \\
\bottomrule
\end{tabular}
\end{table}


The seasonal bias on the DA15 window is striking: BA gives $+26$ EUR/MWh on the critical-flat spread, DDD gives essentially zero; the 2024 calendar-matched placebo year shows the same $+27$ swing in its own pre-vs-post comparison. Almost the entire apparent BA effect is seasonal swing of the spread from summer to winter, not the MTU15 reform. On the IDA15 window the bias is smaller (3--5 EUR/MWh) but goes the opposite direction. In both windows the BA is uninformative about the reform; the DDD isolates the reform from the calendar.

\paragraph{Per-cell OLS-FE and DR-DiD tables.}
\begin{table}[htbp]\centering
\def\sym#1{\ifmmode^{#1}\else\(^{#1}\)\fi}
\caption{OLS-FE DDD: DA clearing price in DA15 window. SE clustered at date.}\label{tab:reform_ddd_da15_da}
\begin{tabular}{l*{3}{cc}}
\toprule
                    &\multicolumn{2}{c}{(1)}     &\multicolumn{2}{c}{(2)}     &\multicolumn{2}{c}{(3)}     \\
                    &\multicolumn{2}{c}{Sparse}  &\multicolumn{2}{c}{Month FE}&\multicolumn{2}{c}{Date FE} \\
                    &           b   &          se&           b   &          se&           b   &          se\\
\midrule
crit*post*y25 ($\beta_7$, DDD)&      -0.286   &     (2.915)&      -0.286   &     (2.916)&      -0.286   &     (2.915)\\
crit*post           &      28.057***&     (2.073)&      28.057***&     (2.074)&      28.057***&     (2.073)\\
crit*y25            &      -3.529*  &     (1.951)&      -3.529*  &     (1.952)&      -3.529*  &     (1.951)\\
post*y25            &     -26.523***&     (5.983)&     -26.523***&     (5.335)&               &            \\
crit                &       0.000   &         (.)&       0.000   &         (.)&       0.000   &         (.)\\
post                &      11.846***&     (4.112)&       0.000   &         (.)&               &            \\
y25                 &      -1.662   &     (4.010)&      -1.662   &     (2.947)&               &            \\
\midrule
Observations        &       13888   &            &       13888   &            &       13888   &            \\
$R^2$               &       0.282   &            &       0.416   &            &       0.800   &            \\
Clock-hour FE       &         Yes   &            &         Yes   &            &         Yes   &            \\
DOW FE              &         Yes   &            &         Yes   &            &         Yes   &            \\
Month FE            &          No   &            &         Yes   &            &         n/a   &            \\
Date FE             &          No   &            &          No   &            &         Yes   &            \\
Cluster             &        Date   &            &        Date   &            &        Date   &            \\
\bottomrule
\end{tabular}
\end{table}

\begin{table}[htbp]\centering
\def\sym#1{\ifmmode^{#1}\else\(^{#1}\)\fi}
\caption{OLS-FE DDD: IDA Session 1 clearing price in DA15 window. SE clustered at date.}\label{tab:reform_ddd_da15_ida}
\begin{tabular}{l*{3}{cc}}
\toprule
                    &\multicolumn{2}{c}{(1)}     &\multicolumn{2}{c}{(2)}     &\multicolumn{2}{c}{(3)}     \\
                    &\multicolumn{2}{c}{Sparse}  &\multicolumn{2}{c}{Month FE}&\multicolumn{2}{c}{Date FE} \\
                    &           b   &          se&           b   &          se&           b   &          se\\
\midrule
crit*post*y25 ($\beta_7$, DDD)&      -2.674   &     (2.852)&      -2.674   &     (2.853)&      -2.674   &     (2.852)\\
crit*post           &      28.703***&     (1.987)&      28.703***&     (1.988)&      28.703***&     (1.987)\\
crit*y25            &      -2.421   &     (1.958)&      -2.421   &     (1.958)&      -2.421   &     (1.958)\\
post*y25            &     -27.432***&     (6.030)&     -26.886***&     (5.313)&               &            \\
crit                &       0.000   &         (.)&       0.000   &         (.)&       0.000   &         (.)\\
post                &      11.870***&     (4.038)&       0.000   &         (.)&               &            \\
y25                 &      -3.158   &     (4.075)&      -3.210   &     (2.961)&               &            \\
\midrule
Observations        &       19138   &            &       19138   &            &       19138   &            \\
$R^2$               &       0.299   &            &       0.479   &            &       0.785   &            \\
Clock-hour FE       &         Yes   &            &         Yes   &            &         Yes   &            \\
DOW FE              &         Yes   &            &         Yes   &            &         Yes   &            \\
Month FE            &          No   &            &         Yes   &            &         n/a   &            \\
Date FE             &          No   &            &          No   &            &         Yes   &            \\
Cluster             &        Date   &            &        Date   &            &        Date   &            \\
\bottomrule
\end{tabular}
\end{table}

\begin{table}[htbp]\centering
\def\sym#1{\ifmmode^{#1}\else\(^{#1}\)\fi}
\caption{OLS-FE DDD: DA clearing price in IDA15 window. SE clustered at date.}\label{tab:reform_ddd_ida15_da}
\begin{tabular}{l*{3}{cc}}
\toprule
                    &\multicolumn{2}{c}{(1)}     &\multicolumn{2}{c}{(2)}     &\multicolumn{2}{c}{(3)}     \\
                    &\multicolumn{2}{c}{Sparse}  &\multicolumn{2}{c}{Month FE}&\multicolumn{2}{c}{Date FE} \\
                    &           b   &          se&           b   &          se&           b   &          se\\
\midrule
crit*post*y25 ($\beta_7$, DDD)&      -2.243   &     (3.630)&      -2.243   &     (3.631)&      -2.243   &     (3.628)\\
crit*post           &     -10.440***&     (1.738)&     -10.440***&     (1.739)&     -10.440***&     (1.738)\\
crit*y25            &       2.657   &     (2.282)&       2.657   &     (2.284)&       2.657   &     (2.282)\\
post*y25            &     -27.426***&     (6.459)&     -27.438***&     (6.250)&               &            \\
crit                &       0.000   &         (.)&       0.000   &         (.)&       0.000   &         (.)\\
post                &     -34.767***&     (4.266)&      -3.141   &     (6.004)&               &            \\
y25                 &      45.676***&     (4.598)&      45.670***&     (4.229)&               &            \\
\midrule
Observations        &        3934   &            &        3934   &            &        3934   &            \\
$R^2$               &       0.573   &            &       0.627   &            &       0.874   &            \\
Clock-hour FE       &         Yes   &            &         Yes   &            &         Yes   &            \\
DOW FE              &         Yes   &            &         Yes   &            &         Yes   &            \\
Month FE            &          No   &            &         Yes   &            &         n/a   &            \\
Date FE             &          No   &            &          No   &            &         Yes   &            \\
Cluster             &        Date   &            &        Date   &            &        Date   &            \\
\bottomrule
\end{tabular}
\end{table}

\begin{table}[htbp]\centering
\def\sym#1{\ifmmode^{#1}\else\(^{#1}\)\fi}
\caption{OLS-FE DDD: IDA Session 1 clearing price in IDA15 window. SE clustered at date.}\label{tab:reform_ddd_ida15_ida}
\begin{tabular}{l*{3}{cc}}
\toprule
                    &\multicolumn{2}{c}{(1)}     &\multicolumn{2}{c}{(2)}     &\multicolumn{2}{c}{(3)}     \\
                    &\multicolumn{2}{c}{Sparse}  &\multicolumn{2}{c}{Month FE}&\multicolumn{2}{c}{Date FE} \\
                    &           b   &          se&           b   &          se&           b   &          se\\
\midrule
crit*post*y25 ($\beta_7$, DDD)&       0.504   &     (3.694)&       0.504   &     (3.695)&       0.504   &     (3.693)\\
crit*post           &      -9.933***&     (1.693)&      -9.933***&     (1.694)&      -9.933***&     (1.693)\\
crit*y25            &       3.266   &     (2.155)&       3.266   &     (2.155)&       3.266   &     (2.154)\\
post*y25            &     -29.831***&     (6.566)&     -29.975***&     (6.342)&               &            \\
crit                &       0.000   &         (.)&       0.000   &         (.)&       0.000   &         (.)\\
post                &     -36.304***&     (4.397)&      -6.630   &     (5.983)&               &            \\
y25                 &      45.197***&     (4.585)&      45.316***&     (4.202)&               &            \\
\midrule
Observations        &        5530   &            &        5530   &            &        5530   &            \\
$R^2$               &       0.588   &            &       0.627   &            &       0.837   &            \\
Clock-hour FE       &         Yes   &            &         Yes   &            &         Yes   &            \\
DOW FE              &         Yes   &            &         Yes   &            &         Yes   &            \\
Month FE            &          No   &            &         Yes   &            &         n/a   &            \\
Date FE             &          No   &            &          No   &            &         Yes   &            \\
Cluster             &        Date   &            &        Date   &            &        Date   &            \\
\bottomrule
\end{tabular}
\end{table}

\begin{table}[H]\centering\small
\caption{DR-DiD on DA clearing price in DA15 window, prices in levels (Sant'Anna--Zhao 2020). 2$\times$2 DiD per hour-class, covariates \texttt{i.dow + i.month}; DDD is the difference of ATTs with joint SE from the influence-function difference.}\label{tab:reform_csdid_da15_da}
\begin{tabular}{lrr}
\toprule
Coefficient & Estimate & SE \\\\
\midrule
ATT$_{\text{crit}}$ &  -8.336 & (4.326) \\\\
ATT$_{\text{flat}}$ &  -2.961 & (4.846) \\\\
DDD = ATT$_{\text{crit}}$ $-$ ATT$_{\text{flat}}$ &  -5.375 & (2.356) \\\\
\midrule
p-value (DDD)        & 0.023 & \\\\
95\% CI (DDD)        & [ -9.992,  -0.758] & \\\\
N (dates)            &   584 & \\\\
Outcome              & DA clearing price & \\\\
Window               & DA15 (reform = 2025-10-01) & \\\\
\bottomrule
\end{tabular}
\end{table}

\begin{table}[H]\centering\small
\caption{DR-DiD on IDA Session 1 clearing price in DA15 window, prices in levels (Sant'Anna--Zhao 2020). 2$\times$2 DiD per hour-class, covariates \texttt{i.dow + i.month}; DDD is the difference of ATTs with joint SE from the influence-function difference.}\label{tab:reform_csdid_da15_ida}
\begin{tabular}{lrr}
\toprule
Coefficient & Estimate & SE \\\\
\midrule
ATT$_{\text{crit}}$ & -10.379 & (4.463) \\\\
ATT$_{\text{flat}}$ &  -5.016 & (4.957) \\\\
DDD = ATT$_{\text{crit}}$ $-$ ATT$_{\text{flat}}$ &  -5.364 & (2.296) \\\\
\midrule
p-value (DDD)        & 0.019 & \\\\
95\% CI (DDD)        & [ -9.864,  -0.863] & \\\\
N (dates)            &   554 & \\\\
Outcome              & IDA Session 1 clearing price & \\\\
Window               & DA15 (reform = 2025-10-01) & \\\\
\bottomrule
\end{tabular}
\end{table}

\begin{table}[H]\centering\small
\caption{DR-DiD on DA clearing price in IDA15 window, prices in levels (Sant'Anna--Zhao 2020). 2$\times$2 DiD per hour-class, covariates \texttt{i.dow + i.month}; DDD is the difference of ATTs with joint SE from the influence-function difference.}\label{tab:reform_csdid_ida15_da}
\begin{tabular}{lrr}
\toprule
Coefficient & Estimate & SE \\\\
\midrule
ATT$_{\text{crit}}$ & -25.019 & (5.486) \\\\
ATT$_{\text{flat}}$ & -31.302 & (6.161) \\\\
DDD = ATT$_{\text{crit}}$ $-$ ATT$_{\text{flat}}$ &   6.283 & (3.009) \\\\
\midrule
p-value (DDD)        & 0.037 & \\\\
95\% CI (DDD)        & [  0.386,  12.180] & \\\\
N (dates)            &   281 & \\\\
Outcome              & DA clearing price & \\\\
Window               & IDA15 (reform = 2025-03-19) & \\\\
\bottomrule
\end{tabular}
\end{table}

\begin{table}[H]\centering\small
\caption{DR-DiD on IDA Session 1 clearing price in IDA15 window, prices in levels (Sant'Anna--Zhao 2020). 2$\times$2 DiD per hour-class, covariates \texttt{i.dow + i.month}; DDD is the difference of ATTs with joint SE from the influence-function difference.}\label{tab:reform_csdid_ida15_ida}
\begin{tabular}{lrr}
\toprule
Coefficient & Estimate & SE \\\\
\midrule
ATT$_{\text{crit}}$ & -25.396 & (5.587) \\\\
ATT$_{\text{flat}}$ & -34.823 & (6.129) \\\\
DDD = ATT$_{\text{crit}}$ $-$ ATT$_{\text{flat}}$ &   9.427 & (2.949) \\\\
\midrule
p-value (DDD)        & 0.001 & \\\\
95\% CI (DDD)        & [  3.647,  15.207] & \\\\
N (dates)            &   278 & \\\\
Outcome              & IDA Session 1 clearing price & \\\\
Window               & IDA15 (reform = 2025-03-19) & \\\\
\bottomrule
\end{tabular}
\end{table}


% =====================================================================
\section{Robustness}\label{sec:ident:price_robustness}
% =====================================================================

\paragraph{Closest-to-delivery IDA price.} The IDA-cell DDDs use Session~1 throughout; Session~1's clearing horizon was redefined at the June~2024 SIDC reform (MIBEL 6-session $\to$ SIDC 3-session). As a robustness, we re-estimate each IDA cell using the maximum-\code{session\_number} price per (date, period). The DA15-IDA spillover cell loses significance ($+1.21$, $p = 0.60$ vs baseline $-5.36^{**}$); the IDA15-IDA direct cell shrinks from $+9.43^{***}$ to $+6.26$ ($p = 0.064$).

\begin{table}[H]\centering\footnotesize
\setlength{\tabcolsep}{4pt}
\caption{DR-DiD on IDA prices: Session 1 (baseline, in Table~\ref{tab:reform_ddd_summary}) vs closest-to-delivery session (robustness). For each (date, period) the closest-to-delivery price is the row with the maximum \code{session_number}: under SIDC, IDA-3 for afternoon and IDA-2 for morning; under MIBEL, Sessions 4--6 for D-day afternoons and Session 3 for morning. This robustness removes the MIBEL-to-SIDC redefinition of Session 1. Covariates: \code{i.dow + i.month}. Standard errors in parentheses. Stars: $^{*}$ 0.10, $^{**}$ 0.05, $^{***}$ 0.01.}
\label{tab:csdid_ida_latest_compare}
\begin{tabular}{lccc}
\toprule
Window & ATT$_{\text{crit}}$ & ATT$_{\text{flat}}$ & DDD \\
\midrule
DA15 &   -2.54 \scriptsize{( 4.39)} &   -3.76 \scriptsize{( 4.75)} &    1.21 \scriptsize{( 2.33)} \\
IDA15 &  -27.84 \scriptsize{( 5.16)} &  -34.10 \scriptsize{( 5.88)} &    6.26* \scriptsize{( 3.38)} \\
\bottomrule
\end{tabular}
\end{table}


The DA-side cells of the headline table are unaffected by this robustness (no session choice in the DA market). We treat the IDA-side numbers as upper bounds until wild-cluster-bootstrap inference and the closest-to-delivery results converge.

\paragraph{Propensity-score overlap.} Across all four cells, $\hat p(X) \in [0.38, 0.54]$ with median $\approx 0.50$; no observations below $0.10$ or above $0.90$. Calendar-matched windows deliver balanced (\text{dow},\,\text{month}) assignment by design; IPW weights are not extreme. Histograms produced as working figures.

\paragraph{Granularity strategy in the panel.} The cross-reform DDD runs on disaggregated period-level data (1 obs/hour pre-MTU15, 4 obs/hour post-MTU15). The 15-min within-hour quarter index is residual heterogeneity absorbed by the FE structure, not a treatment-defining variable. Post-reform within-hour quarter descriptive use (e.g., $D_w$) is a separate complementary exercise from the descriptive companion; it does not enter the DDD.

\paragraph{Next-pass robustness checks (priority order).} (a) Wild-cluster bootstrap (\code{boottest}) on the firm-cluster $\delta_{\text{piv}}$ given $G \approx 9$; (b) reforzada-intensity continuous covariate in the DA15-design DR-DiD outcome model; (c) augmented DR-DiD covariate set with ENTSO-E wind+solar generation and TTF gas as exogenous structural cost shifters; (d) event-study with multiplier-bootstrap simultaneous CI; (e) per-firm decomposition of the pivotal subsample (the partition treats all five Big-4 firms as identical and all four other firms as identical; firm-by-firm is a natural next step).

% =====================================================================
\section{Synthesis: which model predictions are confirmed, wounded, or rejected}\label{sec:ident:synthesis}
% =====================================================================

\begin{description}
\item[Within-hour $D_w$ rises with MTU15-IDA, with the price-setter margin exercised.] Confirmed. $D_w > 0$ rates triple for CCGT critical hours across the three MTU15-IDA regimes, and all four pivotal CCGT firms are within-hour-active under DA15/ID15.
\item[Within-hour bid-shape is concentrated in critical hours.] Confirmed in direction (critical $>$ flat $>$ midday on $D_w$ active rates) but no longer a categorical asymmetry --- flat hours have $D_w > 0$ in 50\% of cells under DA15/ID15. The hour-class taxonomy is validated as an ordering, not as a treated-vs-untreated cut.
\item[$q_2 > 0$ under MTU15: classical Ito-Reguant withholding.] Rejected. $q_2$ is negative for CCGT under both MTU15-IDA and MTU15-DA. The model needs to admit the post-clearing redispatch channel as the outside option, not the IDA leg.
\item[$q_1$ falls for price-setters under MTU15-DA.] Confirmed. The pivotal-status decomposition shows the entire pooled $q_1$ DDD is concentrated in the Big-4 CCGT subsample; non-pivotal firms in the same window show no measurable change.
\item[Critical-flat DA price spread compresses under MTU15-DA.] Confirmed, attributed to the fringe-fill channel rather than to price-setter rent extraction. The compression is consistent with the model's high-granularity prediction only if the fringe is competitive enough to fill the new quarter slots --- which the per-tech CI volume growth (Wind $+36\%$, Solar PV $+108\%$ SA) is consistent with.
\item[Reforzada amplifies the strategic response in ramp-constrained hours (paper~\S2.4 corollary).] Compatible but not separately identified. The DA15-window MTU15-DA DDD bundles the granularity reform with reforzada intensification; the two cannot be cleanly decomposed in the current design. The MTU15-IDA pre-blackout window is the only reforzada-off post-MTU15 sub-window, and it confirms the granularity-only direction of the response on prices and on $q_1/q_2$ before reforzada intensifies.
\end{description}

\paragraph{What the model needs to absorb.} The post-DA Fase~I redispatch channel is empirically the largest pay-off margin for pivotal CCGT firms, and the model as written does not generate it directly --- it appears via the reforzada / RT2 confounder. A model extension that endogenises the outside option (DA participation vs Fase~I redispatch) would deliver the broader-disengagement quadrant as an equilibrium prediction, not as a confound. The CNMC 2019 sanction precedent (descriptive companion, \S\ref{sec:prov:cournot_cnmc:sanctions}) is the regulatory record of this margin being exploited under the older market design; the reforms changed which channel is profitable.

% =====================================================================
\section{Known limitations and the design's remaining open ends}\label{sec:ident:open_ends}
% =====================================================================

\begin{itemize}
\item \emph{Reforzada-intensity drift within the DA15 design.} Reforzada is categorically on in both pre and post DA15 sub-windows but intensifies materially within the design (Fase~I cost $\sim 2.8\times$). The 2024 placebo is a level placebo, not a trend placebo (2024 has no reforzada at all), so the DDD attributes part of the reforzada-intensity drift to MTU15-DA. Addressing this requires a continuous reforzada index as a covariate.
\item \emph{Flat hours are not MTU15-untreated under DA15/ID15.} CCGT flat $D_w > 0$ is $50.6\%$ vs critical $70.7\%$. The DA15 DDD identifies a critical-flat differential between two treated arms, not the absolute critical-hour effect. The level of the absolute critical-hour effect is bounded below by the DDD coefficient and above by the BA-with-controls estimate.
\item \emph{Pivotal-status heterogeneity has lost some bite.} All four pivotal CCGT firms now show strong within-hour activity under DA15/ID15; the cross-firm response is more uniform than the earlier ``GN only'' reading suggested. The pivotal-coefficient is close to the pooled coefficient, and the non-pivotal-as-internal-control reading is consequently less informative for the granularity-exploitation channel.
\item \emph{Short MTU15-IDA pre-blk post-window.} 40 days. Many specs hit power limits; wild-cluster bootstrap is the planned robustness.
\item \emph{Spillover cells are interpretive, not headline.} The DA15-IDA and IDA15-DA cells test cross-market spillover and are reported in the headline table for completeness; their interpretation depends on a forward-looking-bidding assumption that we do not formally model.
\end{itemize}

% =====================================================================
\section{Supporting tables (legacy version)}\label{sec:ident:legacy}
% =====================================================================

\noindent A more discursive presentation of the same DDD pipeline lives in \code{thesis/provisional/regression_results_legacy.tex}. The legacy file is preserved for the commit-history record; the table outputs cited above are produced by the same Stata pipeline (\code{scripts/stata/reform_window/}).

\end{document}
